% 致谢
\begin{acknowledgement}
% 【待填充 / 可自行润色】

时光荏苒, 在上海大学机电工程与自动化学院攻读硕士学位的时光即将结束. 回首这段求
学历程, 感慨良多, 心中充满了感激之情.

首先, 我要向我的导师王曰英教授致以最诚挚的感谢. 王老师严谨的治学态度、渊博的学
识以及对学术前沿的敏锐洞察深深地影响了我. 从课题方向的确立、论文框架的搭建, 到
仿真实验与物理实验的具体细节, 王老师都给予了我耐心细致的指导. 尤为令我难忘的
是, 在我面对研究瓶颈时, 王老师总能以独到的见解帮我指明方向; 在论文投稿遇到挫折
时, 王老师的鼓励给了我坚持下去的动力. 师恩难忘, 谨此致以最深切的谢意.

其次, 感谢课题组各位老师在学习与科研过程中的无私帮助. 感谢实验室同门师兄弟姐妹们
在日常讨论、算法调试与实验开展中的通力合作, 在与你们的交流与思想碰撞中, 我的学术
视野得到了极大拓展.

再次, 感谢参与本文多无人艇物理实验平台搭建与湖面实验的团队成员们. 从无人艇硬件
选型、控制系统调试, 到现场实验, 每一步都凝聚了大家的心血. 同时, 也要感谢为本研究
提供资助的国家自然科学基金优秀青年基金项目 (62122046).

最后, 我要向我的家人致以最深的感谢. 父母的默默支持与理解是我求学道路上最坚实
的后盾; 家人的关心与包容让我能够心无旁骛地投身学术. 你们的爱是我前行的最大动力.

谨以此文献给所有关心、帮助和支持过我的人. 在未来的科研与工作道路上, 我将继续秉
承上海大学 ``自强不息, 先天下之忧而忧'' 的校训, 不忘初心, 砥砺前行.

  \vspace{50bp}

  \hfill 胡恒宇

  \hfill 上海大学 机电工程与自动化学院

  \hfill 2026 年 6 月

\end{acknowledgement}
